
This paper addresses the absence of structural diagnostic tools for comparing code generation alignment methods. It introduces structural efficiency — semantic AST edit distance per unit KL divergence — as a measurement quantity, and implements a five-module framework (FA-AST classification, ZSS edit distance, KL-matched checkpoint selection, SEP analysis, bootstrap and Mann-Whitney testing) that has been validated end-to-end on real DeepSeek-Coder-7B-instruct-v1.5 checkpoints.

Two sub-experiments were executed. The proof-of-concept (h-e1) confirmed that the measurement pipeline correctly detects structural differences when they are engineered into synthetic inputs, yielding a bootstrap 95\% CI of [4.65, 8.73] for the mean edit-per-KL differential. The mechanism analysis (h-m1) produced SEP ≈ 0.237 for both GRPO and DPO on real checkpoints, with p = 0.4248 (Mann-Whitney), but was found to be underpowered due to checkpoint aliasing that reduced the effective sample size from 27 nominal pairs to approximately 2. The mechanism hypothesis — that execution reward selectively concentrates policy movement on control-flow and data-flow AST transformations — remains empirically untested.

The paper additionally documents checkpoint aliasing as a previously undocumented confound in RL fine-tuning checkpoint-comparison studies and provides a one-assertion safeguard that prevents silent corruption of nominally large analyses.

\textbf{Future directions.} The immediate priority is a corrected experimental run: 1000-step GRPO and DPO training with \texttt{save_steps=100}, real evalplus-based execution rewards, and execution-oracle DPO pairs, followed by h-m1 analysis on 10+ unique checkpoints. If the corrected run confirms SEP_GRPO > SEP_DPO with adequate statistical power, the downstream mechanism chain (SEP as pass@1 mediator, ECE correlation, OOD transfer on LiveCodeBench) becomes experimentally tractable using the same pipeline infrastructure. If near-equal SEP persists with adequate power, this constitutes a genuine finding that GRPO's structural advantage lies in total edit volume rather than selective semantic node concentration.

The framework is also applicable beyond GRPO versus DPO: any two post-training methods producing AST-parseable code checkpoints can be compared. Instruction tuning, PPO, rejection sampling fine-tuning, and DPO variants are all amenable to structural efficiency analysis using the same pipeline.
